% Options for packages loaded elsewhere
% Options for packages loaded elsewhere
\PassOptionsToPackage{unicode}{hyperref}
\PassOptionsToPackage{hyphens}{url}
\PassOptionsToPackage{dvipsnames,svgnames,x11names}{xcolor}
%
\documentclass[
  letterpaper,
  DIV=11,
  numbers=noendperiod]{scrartcl}
\usepackage{xcolor}
\usepackage{amsmath,amssymb}
\setcounter{secnumdepth}{-\maxdimen} % remove section numbering
\usepackage{iftex}
\ifPDFTeX
  \usepackage[T1]{fontenc}
  \usepackage[utf8]{inputenc}
  \usepackage{textcomp} % provide euro and other symbols
\else % if luatex or xetex
  \usepackage{unicode-math} % this also loads fontspec
  \defaultfontfeatures{Scale=MatchLowercase}
  \defaultfontfeatures[\rmfamily]{Ligatures=TeX,Scale=1}
\fi
\usepackage{lmodern}
\ifPDFTeX\else
  % xetex/luatex font selection
\fi
% Use upquote if available, for straight quotes in verbatim environments
\IfFileExists{upquote.sty}{\usepackage{upquote}}{}
\IfFileExists{microtype.sty}{% use microtype if available
  \usepackage[]{microtype}
  \UseMicrotypeSet[protrusion]{basicmath} % disable protrusion for tt fonts
}{}
\makeatletter
\@ifundefined{KOMAClassName}{% if non-KOMA class
  \IfFileExists{parskip.sty}{%
    \usepackage{parskip}
  }{% else
    \setlength{\parindent}{0pt}
    \setlength{\parskip}{6pt plus 2pt minus 1pt}}
}{% if KOMA class
  \KOMAoptions{parskip=half}}
\makeatother
% Make \paragraph and \subparagraph free-standing
\makeatletter
\ifx\paragraph\undefined\else
  \let\oldparagraph\paragraph
  \renewcommand{\paragraph}{
    \@ifstar
      \xxxParagraphStar
      \xxxParagraphNoStar
  }
  \newcommand{\xxxParagraphStar}[1]{\oldparagraph*{#1}\mbox{}}
  \newcommand{\xxxParagraphNoStar}[1]{\oldparagraph{#1}\mbox{}}
\fi
\ifx\subparagraph\undefined\else
  \let\oldsubparagraph\subparagraph
  \renewcommand{\subparagraph}{
    \@ifstar
      \xxxSubParagraphStar
      \xxxSubParagraphNoStar
  }
  \newcommand{\xxxSubParagraphStar}[1]{\oldsubparagraph*{#1}\mbox{}}
  \newcommand{\xxxSubParagraphNoStar}[1]{\oldsubparagraph{#1}\mbox{}}
\fi
\makeatother


\usepackage{longtable,booktabs,array}
\usepackage{calc} % for calculating minipage widths
% Correct order of tables after \paragraph or \subparagraph
\usepackage{etoolbox}
\makeatletter
\patchcmd\longtable{\par}{\if@noskipsec\mbox{}\fi\par}{}{}
\makeatother
% Allow footnotes in longtable head/foot
\IfFileExists{footnotehyper.sty}{\usepackage{footnotehyper}}{\usepackage{footnote}}
\makesavenoteenv{longtable}
\usepackage{graphicx}
\makeatletter
\newsavebox\pandoc@box
\newcommand*\pandocbounded[1]{% scales image to fit in text height/width
  \sbox\pandoc@box{#1}%
  \Gscale@div\@tempa{\textheight}{\dimexpr\ht\pandoc@box+\dp\pandoc@box\relax}%
  \Gscale@div\@tempb{\linewidth}{\wd\pandoc@box}%
  \ifdim\@tempb\p@<\@tempa\p@\let\@tempa\@tempb\fi% select the smaller of both
  \ifdim\@tempa\p@<\p@\scalebox{\@tempa}{\usebox\pandoc@box}%
  \else\usebox{\pandoc@box}%
  \fi%
}
% Set default figure placement to htbp
\def\fps@figure{htbp}
\makeatother





\setlength{\emergencystretch}{3em} % prevent overfull lines

\providecommand{\tightlist}{%
  \setlength{\itemsep}{0pt}\setlength{\parskip}{0pt}}



 


\KOMAoption{captions}{tableheading}
\usepackage{tikz}
\usetikzlibrary{arrows.meta,calc,intersections}
\makeatletter
\@ifpackageloaded{caption}{}{\usepackage{caption}}
\AtBeginDocument{%
\ifdefined\contentsname
  \renewcommand*\contentsname{Table of contents}
\else
  \newcommand\contentsname{Table of contents}
\fi
\ifdefined\listfigurename
  \renewcommand*\listfigurename{List of Figures}
\else
  \newcommand\listfigurename{List of Figures}
\fi
\ifdefined\listtablename
  \renewcommand*\listtablename{List of Tables}
\else
  \newcommand\listtablename{List of Tables}
\fi
\ifdefined\figurename
  \renewcommand*\figurename{Figure}
\else
  \newcommand\figurename{Figure}
\fi
\ifdefined\tablename
  \renewcommand*\tablename{Table}
\else
  \newcommand\tablename{Table}
\fi
}
\@ifpackageloaded{float}{}{\usepackage{float}}
\floatstyle{ruled}
\@ifundefined{c@chapter}{\newfloat{codelisting}{h}{lop}}{\newfloat{codelisting}{h}{lop}[chapter]}
\floatname{codelisting}{Listing}
\newcommand*\listoflistings{\listof{codelisting}{List of Listings}}
\makeatother
\makeatletter
\makeatother
\makeatletter
\@ifpackageloaded{caption}{}{\usepackage{caption}}
\@ifpackageloaded{subcaption}{}{\usepackage{subcaption}}
\makeatother
\usepackage{bookmark}
\IfFileExists{xurl.sty}{\usepackage{xurl}}{} % add URL line breaks if available
\urlstyle{same}
\hypersetup{
  pdftitle={AREC 310 - Homework 3 - Answer Key},
  pdfauthor={Lauren Chenarides},
  colorlinks=true,
  linkcolor={blue},
  filecolor={Maroon},
  citecolor={Blue},
  urlcolor={Blue},
  pdfcreator={LaTeX via pandoc}}


\title{AREC 310 - Homework 3 - Answer Key}
\usepackage{etoolbox}
\makeatletter
\providecommand{\subtitle}[1]{% add subtitle to \maketitle
  \apptocmd{\@title}{\par {\large #1 \par}}{}{}
}
\makeatother
\subtitle{30 points possible}
\author{}
\date{}
\begin{document}
\maketitle


Directions: Answer the following questions completely. You can work with
your peers, but you MUST TURN IN YOUR OWN ANSWERS (i.e., in your own
words) to get credit. Upload a PDF to Canvas by the due date.

\subsection{Problem 1: Demand (5
points)}\label{problem-1-demand-5-points}

You are given the following demand schedule for white, whole grain
processed flour. Use it to answer the following questions.

\begin{longtable}[]{@{}ll@{}}
\toprule\noalign{}
Price (U.S. dollars per lb) & Quantity (1000 lbs) \\
\midrule\noalign{}
\endhead
\bottomrule\noalign{}
\endlastfoot
15 & 10 \\
16 & 8 \\
17 & 6 \\
18 & 4 \\
19 & 2 \\
20 & 0 \\
\end{longtable}

\begin{enumerate}
\def\labelenumi{\alph{enumi}.}
\tightlist
\item
  Plot the demand curve using appropriate economic conventions.
\end{enumerate}

Demand is linear from \((Q=10, P=15)\) to \((Q=0, P=20)\).

\begin{figure}[h]
\centering
\begin{tikzpicture}[>=Latex, scale=0.9]
  % axes
  \draw[->] (0,0) -- (11,0) node[below] {$Q$ (thousand lbs)};
  \draw[->] (0,0) -- (0,8) node[left] {$P$ (\$/lb)};
  % ticks
  \foreach \x/\lab in {0/0,2/2,4/4,6/6,8/8,10/10} \draw (\x,0) -- ++(0,0.1) node[below,yshift=-2pt]{\lab};
  \foreach \y/\lab in {0/0,3/15,4/18,5/20} \draw (0,\y) -- ++(0.1,0) node[left,xshift=-8pt]{\lab};
  % demand line
  \draw[thick,blue] (0,5) -- (10,3);
  % labeled endpoints
  \fill (10,3) circle (1.5pt) node[below right] {$Q{=}10,\ P{=}15$};
  \fill (0,5) circle (1.5pt) node[above left] {$Q{=}0,\ P{=}20$};
  \node[blue] at (7.5,3.4) {$D$};
\end{tikzpicture}
\caption{Demand for white whole grain processed flour.}
\end{figure}

\begin{enumerate}
\def\labelenumi{\alph{enumi}.}
\setcounter{enumi}{1}
\tightlist
\item
  What is the slope? Is it positive or negative? Why?
\end{enumerate}

With (P) on the vertical axis and (Q) on the horizontal axis, slope is
(\Delta P/\Delta Q = (20-15)/(0-10) = -0.5). Negative because of the law
of demand (price and quantity demanded move in opposite directions).

\begin{enumerate}
\def\labelenumi{\alph{enumi}.}
\setcounter{enumi}{2}
\tightlist
\item
  Calculate the own-price elasticity of demand when the price changes
  from \$18/lb to \$15/lb. Is it elastic or inelastic? Interpret.
\end{enumerate}

From table: at (P\_1=18), (Q\_1=4); at (P\_2=15), (Q\_2=10).

{[} E\_D=\frac{\frac{Q_2-Q_1}{(Q_1+Q_2)/2}}{\frac{P_2-P_1}{(P_1+P_2)/2}}
=\frac{(10-4)/7}{(15-18)/16.5}\approx \frac{0.8571}{-0.1818}\approx -4.71
{]} Elastic (absolute value (\textgreater1)). Interpretation: a 1\% drop
in price increases quantity demanded by about 4.7\%.

\begin{enumerate}
\def\labelenumi{\alph{enumi}.}
\setcounter{enumi}{3}
\tightlist
\item
  ConAgra (the owner of Ultragrain, a version of processed whole wheat
  flour) recently began an advertising campaign highlighting the health
  benefits of increasing whole grains in a diet. Will this affect the
  demand schedule given above? If so, how? Clearly illustrate on your
  graph from (a). If not, explain why.
\end{enumerate}

Shifts \textbf{demand right} (higher quantity at each price). On the
graph: a parallel rightward shift of the demand curve.

\begin{enumerate}
\def\labelenumi{\alph{enumi}.}
\setcounter{enumi}{4}
\tightlist
\item
  Researchers at CSU recently developed a new method to clean the
  haulers that transport the white, whole wheat grain from the farms to
  the processing plants that decreases the amount of time it takes to
  clean the haulers. Will this affect the demand schedule given above?
  If so, how? Clearly illustrate on your graph from (a). If not, explain
  why.
\end{enumerate}

This is a \textbf{supply-side} cost/time change, not a demand
determinant. Demand schedule \textbf{unchanged}.

\subsection{Problem 2: Supply (5
points)}\label{problem-2-supply-5-points}

You are given the following schedule of supply for white, whole grain
processed flour. Use it to answer the following questions.

\begin{longtable}[]{@{}ll@{}}
\toprule\noalign{}
Price (U.S. dollars per lb) & Quantity (1000 lbs) \\
\midrule\noalign{}
\endhead
\bottomrule\noalign{}
\endlastfoot
34 & 12 \\
30 & 10 \\
26 & 8 \\
22 & 6 \\
18 & 4 \\
14 & 2 \\
\end{longtable}

\begin{enumerate}
\def\labelenumi{\alph{enumi}.}
\tightlist
\item
  Plot the supply curve using appropriate economic conventions.
\end{enumerate}

Supply is linear from \((Q=2, P=14)\) to \((Q=12, P=34)\).

\begin{figure}[h]
\centering
\begin{tikzpicture}[>=Latex, scale=0.9]
  % axes
  \draw[->] (0,0) -- (13,0) node[below] {$Q$ (thousand lbs)};
  \draw[->] (0,0) -- (0,10) node[left] {$P$ (\$/lb)};
  % ticks
  \foreach \x/\lab in {0/0,2/2,4/4,6/6,8/8,10/10,12/12} \draw (\x,0) -- ++(0,0.1) node[below,yshift=-2pt]{\lab};
  \foreach \y/\lab in {0/0,3/14,6/22,8.5/30,10/34} \draw (0,\y) -- ++(0.1,0) node[left,xshift=-8pt]{\lab};
  % supply line
  \draw[thick,green!50!black] (2,3) -- (12,10);
  % labeled endpoints
  \fill (2,3) circle (1.5pt) node[left] {$Q{=}2,\ P{=}14$};
  \fill (12,10) circle (1.5pt) node[above right] {$Q{=}12,\ P{=}34$};
  \node[green!50!black] at (9,8.2) {$S$};
\end{tikzpicture}
\caption{Supply for white whole grain processed flour.}
\end{figure}

\begin{enumerate}
\def\labelenumi{\alph{enumi}.}
\setcounter{enumi}{1}
\tightlist
\item
  What is the slope? Is it positive or negative? Why?
\end{enumerate}

With (P) vertical, (Q) horizontal: (\Delta P/\Delta Q = (34-14)/(12-2) =
20/10 = 2) (positive, law of supply).

\begin{enumerate}
\def\labelenumi{\alph{enumi}.}
\setcounter{enumi}{2}
\tightlist
\item
  Calculate the own-price elasticity of supply when the price changes
  from \$22/lb to \$30/lb. Is it elastic or inelastic? Interpret.
\end{enumerate}

From table: (P\_1=22, Q\_1=6); (P\_2=30, Q\_2=10).

{[} E\_S=\frac{\frac{10-6}{(6+10)/2}}{\frac{30-22}{(22+30)/2}}
=\frac{4/8}{8/26}\approx \frac{0.5}{0.3077}\approx 1.63 {]} Elastic
((\textgreater1)). Interpretation: a 1\% rise in price raises quantity
supplied by about 1.6\%.

\begin{enumerate}
\def\labelenumi{\alph{enumi}.}
\setcounter{enumi}{3}
\tightlist
\item
  ConAgra recently began an advertising campaign highlighting the health
  benefits of increasing whole grains in a diet. Will this affect the
  supply schedule given above? If so, how? Clearly illustrate on your
  graph from (a). If not, explain why.
\end{enumerate}

Advertising is a \textbf{demand} determinant. The \textbf{supply}
schedule does \textbf{not} shift.

\begin{enumerate}
\def\labelenumi{\alph{enumi}.}
\setcounter{enumi}{4}
\tightlist
\item
  Researchers at CSU recently developed a new method to clean the
  haulers that transport the white, whole wheat grain from the farms to
  the processing plants that decreases the amount of time it takes to
  clean the haulers. Will this affect the supply schedule given above?
  If so, how? Clearly illustrate on your graph from (a). If not, explain
  why.
\end{enumerate}

Lowers cost/time per unit → \textbf{supply shifts right} (greater
quantity at each price).

\subsection{Problem 3: Equilibrium (9
points)}\label{problem-3-equilibrium-9-points}

Given the supply and demand conditions illustrated in Problems 1 and 2,
respectively, answer the following questions.

Let (Q\_D) and (Q\_S) be (thousand lbs), (P) in \$/lb.

From Problem 1 (linear), (dQ/dP=-2) and point ((P,Q)=(15,10)) ⇒ (Q\_D =
40 - 2P).

From Problem 2 (linear), (dQ/dP=+0.5) and point ((14,2)) ⇒ (Q\_S = -5 +
0.5P).

\begin{enumerate}
\def\labelenumi{\alph{enumi}.}
\tightlist
\item
  What is the equilibrium price and quantity in the white, whole wheat
  processed flour market under the initial conditions depicted by the
  demand and supply schedules? Graph it using the complete schedules.
\end{enumerate}

Solve (Q\_D = Q\_S): {[} 40 - 2P = -5 + 0.5P ;\Rightarrow; 45 = 2.5P
;\Rightarrow; P\^{}* = 18 {]} {[} Q\^{}* = -5 + 0.5(18) = 4 {]}

\begin{figure}[h]
\centering
\begin{tikzpicture}[>=Latex, scale=0.95]
  % axes
  \draw[->] (0,0) -- (11,0) node[below] {$Q$ (thousand lbs)};
  \draw[->] (0,0) -- (0,8) node[left] {$P$ (\$/lb)};
  % ticks
  \foreach \x in {0,2,4,6,8,10} \draw (\x,0) -- ++(0,0.1) node[below,yshift=-2pt]{\x};
  \foreach \y/\lab in {0/0,3/15,4/18,5/20,6/22,8/26} \draw (0,\y) -- ++(0.1,0) node[left,xshift=-8pt]{\lab};
  % demand (Q=40-2P -> P=(40-Q)/2)
  \draw[thick,blue] (0,5) -- (10,3);
  % supply (Q=-5+0.5P -> P=2(Q+5))
  \draw[thick,green!50!black] (0,10) -- (10,30); % off-scale line; instead plot between known points:
  \draw[thick,green!50!black] (2,3) -- (12,10);  % main visible segment
  % equilibrium point (Q=4,P=18)
  \fill (4,4) circle (2pt) node[above right] {$E(Q^*=4,\ P^*=18)$};
  % guides
  \draw[densely dotted] (4,0) -- (4,4);
  \draw[densely dotted] (0,4) -- (4,4);
  % labels
  \node[blue] at (8,3.4) {$D$};
  \node[green!50!black] at (9.8,8.6) {$S$};
\end{tikzpicture}
\caption{Equilibrium at \(P^*=18\), \(Q^*=4\).}
\end{figure}

\begin{enumerate}
\def\labelenumi{\alph{enumi}.}
\setcounter{enumi}{1}
\tightlist
\item
  Assume white whole wheat processed flour is a normal good. If food
  assistance benefits increase, will the equilibrium in the market be
  different from what you found in part (a)? If so, how? Clearly
  illustrate on your graph. If not, explain why.
\end{enumerate}

Demand shifts \textbf{right} → (P\uparrow, Q\uparrow). Show a parallel
outward shift of (D) on the graph, new intersection above/right of (E).

\begin{enumerate}
\def\labelenumi{\alph{enumi}.}
\setcounter{enumi}{2}
\tightlist
\item
  Assume that more farmers start to plant white, whole wheat. Assume
  that the planting and harvesting conditions were average. Will the
  equilibrium in the market be different from what you found in part
  (a)? If so, how? Clearly illustrate on your graph. If not, explain
  why.
\end{enumerate}

Supply shifts \textbf{right} → (P\downarrow, Q\uparrow). Show a parallel
outward shift of (S), new intersection below/right of (E).

\subsection{Problem 4: Impact of Higher Beef Prices on the Pork Market
(11
points)}\label{problem-4-impact-of-higher-beef-prices-on-the-pork-market-11-points}

Suppose you are a pork market analyst, and beef prices are expected to
rise dramatically, by 10\%.

Because beef and pork are substitutes, this price change creates a
demand shock in the pork market. Specifically, you are asked to
calculate the resulting changes in the pork price and pork quantity.

\textbf{Parameters:}

\begin{itemize}
\tightlist
\item
  Long-run supply elasticity of pork: \(E_S = 2.15\)\\
\item
  Long-run demand elasticity of pork: \(E_D = -1.96\)\\
\item
  Cross-price elasticity of pork demand with respect to beef: \(0.6\)
\end{itemize}

\textbf{Tasks:}

\begin{enumerate}
\def\labelenumi{\arabic{enumi}.}
\tightlist
\item
  Calculate the value of the demand shock, \(S_D\).\\
\item
  Write the supply and demand equations in elasticity form.\\
\item
  Solve for the equilibrium percent change in pork price,
  \(\% \Delta P\).\\
\item
  Solve for the equilibrium percent change in pork quantity,
  \(\% \Delta Q\).\\
\item
  Briefly explain the intuition behind your results.
\end{enumerate}

\emph{Hint: Recall that
\(S_D = (\text{cross-price elasticity})(\% \Delta P_{\text{beef}})\)}

\subsection{Answer Key}\label{answer-key}

\subsubsection{Step 1. Demand shock}\label{step-1.-demand-shock}

\[
S_D = (0.6)(10) = 6 \%
\]

Because the shock comes from beef prices (a substitute), it shifts the
pork demand curve to the right.\\
Supply is unaffected, so \(S_S = 0\).

\subsubsection{Step 2. Write supply and demand
equations}\label{step-2.-write-supply-and-demand-equations}

Supply:

\[
\% \Delta Q_S = E_S (\% \Delta P) + S_S = 2.15 (\% \Delta P) + 0
\]

Demand:

\[
\% \Delta Q_D = E_D (\% \Delta P) + S_D = -1.96 (\% \Delta P) + 6
\]

\subsubsection{\texorpdfstring{Step 3. Impose equilibrium and solve for
\(\% \Delta P\)}{Step 3. Impose equilibrium and solve for \textbackslash\% \textbackslash Delta P}}\label{step-3.-impose-equilibrium-and-solve-for-delta-p}

\[
\begin{aligned}
\% \Delta Q_S &= \% \Delta Q_D \\
2.15 (\% \Delta P) &= -1.96 (\% \Delta P) + 6 \\
(2.15 + 1.96)(\% \Delta P) &= 6 \\
\% \Delta P &= \frac{6}{4.11} \\
\% \Delta P &\approx 1.46 \%
\end{aligned}
\]

So pork prices are expected to rise \textbf{1.46\%}.

\subsubsection{\texorpdfstring{Step 4. Solve for
\(\% \Delta Q\)}{Step 4. Solve for \textbackslash\% \textbackslash Delta Q}}\label{step-4.-solve-for-delta-q}

Supply side:

\[
\% \Delta Q_S = 2.15 (1.46) = 3.14 \%
\]

Demand side:

\[
\% \Delta Q_D = -1.96(1.46) + 6 = -2.86 + 6 = 3.14 \%
\]

Both agree: quantity rises by \textbf{3.14\%}.

\subsubsection{Step 5. Intuition}\label{step-5.-intuition}

\begin{itemize}
\tightlist
\item
  The 10\% increase in beef prices makes pork more attractive.\\
\item
  Demand shifts rightward (\(S_D = +6\%\)).\\
\item
  Pork prices rise modestly (+1.46\%), while pork output increases more
  strongly (+3.14\%).\\
\item
  The results are consistent with elasticities: because supply is
  relatively elastic in the long run, much of the adjustment occurs
  through higher output rather than sharply higher prices.
\end{itemize}




\end{document}
